% ---------------------------------------------------------------------------
% Preâmbulo: pacotes e macros compartilhados por parcial.tex e final.tex
% Compilado com XeLaTeX (ver latexmkrc), por causa da fonte Arial.
% ---------------------------------------------------------------------------

% ---------------------------------------------------------------------------
% Fontes. O texto é Arial 12 (o tamanho vem da opção 12pt da classe). Numa
% máquina sem Arial instalada, usa-se a TeX Gyre Heros, clone livre da
% Helvetica, de métrica equivalente. A capa mantém a fonte serifada com que foi
% desenhada (TeX Gyre Termes, clone livre da Times).
% ---------------------------------------------------------------------------
\usepackage{amsmath}
\usepackage{fontspec}
\IfFontExistsTF{Arial}{%
  \setmainfont{Arial}
  \setsansfont{Arial}
}{%
  \setmainfont{texgyreheros}[Extension=.otf, UprightFont=*-regular,
    BoldFont=*-bold, ItalicFont=*-italic, BoldItalicFont=*-bolditalic]
  \setsansfont{texgyreheros}[Extension=.otf, UprightFont=*-regular,
    BoldFont=*-bold, ItalicFont=*-italic, BoldItalicFont=*-bolditalic]
}
\newfontfamily\fontecapa{texgyretermes}[Extension=.otf, UprightFont=*-regular,
  BoldFont=*-bold, ItalicFont=*-italic, BoldItalicFont=*-bolditalic]
% Letras e dígitos das poucas fórmulas do texto na mesma fonte do texto
\usepackage[italic]{mathastext}

\usepackage{indentfirst}
% Folga para parágrafos com fórmulas ou termos longos que não se quebram
\setlength{\emergencystretch}{2em}
\usepackage{microtype}
\usepackage{graphicx}
\usepackage{xcolor}
\usepackage{booktabs}
\usepackage{tabularx}
% Coluna de parágrafo alinhada à esquerda, para tabelas com texto (evita a
% hifenização forçada das colunas estreitas justificadas)
\newcolumntype{L}[1]{>{\raggedright\arraybackslash}p{#1}}
\usepackage{listings}
\usepackage{tikz}
\usetikzlibrary{positioning,arrows.meta,fit,backgrounds}

% Citações no padrão ABNT, autor-data
\usepackage[brazilian,hyperpageref]{backref}
\usepackage[alf]{abntex2cite}

\renewcommand{\backrefpagesname}{Citado na(s) página(s):~}
\renewcommand{\backref}{}
\renewcommand*{\backrefalt}[4]{%
  \ifcase #1 %
    Nenhuma citação no texto.%
  \or
    Citado na página #2.%
  \else
    Citado #1 vezes nas páginas #2.%
  \fi}

% ---------------------------------------------------------------------------
% Estética: fora da capa, tudo em preto. As duas cores abaixo são usadas
% apenas pela capa (config/capa.tex).
% ---------------------------------------------------------------------------
\definecolor{azulpoli}{HTML}{0B2A4A}
\definecolor{cinzacapa}{HTML}{5A6270}

\renewcommand{\ABNTEXchapterfont}{\bfseries}
\renewcommand{\ABNTEXsectionfont}{\bfseries}
\renewcommand{\ABNTEXsubsectionfont}{\bfseries}

\AtBeginDocument{%
  \hypersetup{
    pdftitle={\imprimirtitulo},
    pdfauthor={\autorumnome, \autordoisnome},
    pdfsubject={\imprimirtipotrabalho},
    hidelinks,
  }%
}

% Código C nas listagens: sem cor e sem negrito
\lstset{
  language=C,
  basicstyle=\ttfamily\small,
  keywordstyle=,
  commentstyle=\itshape,
  numbers=left,
  numberstyle=\tiny,
  frame=single,
  breaklines=true,
  showstringspaces=false,
}

% ---------------------------------------------------------------------------
% Marcação de trabalho pendente.
% Para a versão de entrega, troque \mostrarpendentestrue por \mostrarpendentesfalse:
% os marcadores somem do PDF sem precisar apagá-los do texto.
% ---------------------------------------------------------------------------
\newif\ifmostrarpendentes
\mostrarpendentesfalse
\newcommand{\pendente}[1]{%
  \ifmostrarpendentes{\color{red}[PENDENTE: #1]}\fi}

% Sem negrito no corpo do texto: os rótulos de listas de descrição (QP1, nomes
% das métricas) saem em estilo normal, seguidos de dois-pontos.
\renewcommand*{\descriptionlabel}[1]{\hspace{\labelsep}\normalfont #1:}

% Termos recorrentes
\newcommand{\isr}{\textit{ISR}}
\newcommand{\llm}{\textit{LLM}}
